\section{Conclusioni}
L'esperienza realizzata ha permesso di verificare con buoni livelli di confidenza diverse leggi dell'ottica geometrica. Sono state verificate sperimentalmente la legge di Snell e alcune conseguenze dirette, ovvero la reversibilità del cammino ottico, la riflessione totale e la deviazione di un raggio incidente in un prisma a facce parallele. Dalla legge di Snell e dal metodo dell'angolo di deviazione minima sono state ottenute con buona accuratezza stime indipendenti degli indici di rifrazione del plexiglass e del vetro. Le misure realizzate su due lenti sottili in plastica hanno confermato la validità della legge dei punti coniugati. I valori ottenuti per gli ingrandimenti e per la distanza minima si sono dimostrati statisticamente compatibili con i modelli teorici.
\newpage